% GENERATED FILE. Edit canonical lessons, then run npm run generate:books.
% canonical-source-hash: fnv1a64:6822ebc8e9b7e3bb
% canonical-lessons: RU-C147-bolnoy, RU-C147-prostuda, RU-C147-gripp, RU-C147-nasmork, RU-C147-tabletka

\chapter{Ill, Cold, Flu, Runny nose, Pill}
\label{ch:ru-ill-cold-flu-runny-nose-pill}
\hlchaptermodality{147}

\begin{chapteropening}
\textbf{By the end of this chapter:} \emph{I can say ill, a cold, flu, a runny nose and a pill.}

Complete the last lesson of chapter 147: I can say ill, a cold, flu, a runny nose and a pill.
\end{chapteropening}

\section[\texorpdfstring{\ru{больной} (bolnóy)}{больной (bolnóy)}]{\ru{больной} (bolnóy) — ill; a patient}

\label{lesson:RU-C147-bolnoy}



\begin{quote}
Before the new one: say the Russian for to stop, then the Russian for to park.
\end{quote}

\subsection*{You'll want to know: \ru{больной}}
\textbf{\ru{больной}} (\emph{bolnóy}) — \textquotedblleft{}ill; a patient\textquotedblright{}.

\begin{grammarlens}[title={What you've built}]
One more word for this chapter.
\end{grammarlens}

\subsection*{Your turn}
\begin{itemize}
\raggedright
  \item \emph{Say it:} \emph{bolnóy}
  \item \emph{Say it:} \emph{bolnóy}, once more
  \item \emph{Say it:} say \emph{bolnóy}
  \item \emph{Recall:} say the Russian for to stop, then the Russian for to park, then say \emph{bolnóy} again
\end{itemize}

\begin{tcolorbox}[breakable,colback=teal!4,colframe=teal!35!black,title={Before you move on}]
What does \ru{больной} mean? (\textquotedblleft{}Ill; a patient\textquotedblright{}.) Say it once more.
\end{tcolorbox}
\section[\texorpdfstring{\ru{простуда} (prostúda)}{простуда (prostúda)}]{\ru{простуда} (prostúda) — a cold}

\label{lesson:RU-C147-prostuda}



\begin{quote}
Before the new one: say the Russian for to park, then the Russian for ill; a patient.
\end{quote}

\subsection*{You'll want to know: \ru{простуда}}
\textbf{\ru{простуда}} (\emph{prostúda}) — \textquotedblleft{}a cold\textquotedblright{}.

\begin{grammarlens}[title={What you've built}]
One more word for this chapter.
\end{grammarlens}

\subsection*{Your turn}
\begin{itemize}
\raggedright
  \item \emph{Say it:} \emph{prostúda}
  \item \emph{Say it:} \emph{prostúda}, once more
  \item \emph{Say it:} say \emph{prostúda}
  \item \emph{Recall:} say the Russian for to park, then the Russian for ill; a patient, then say \emph{prostúda} again
\end{itemize}

\begin{tcolorbox}[breakable,colback=teal!4,colframe=teal!35!black,title={Before you move on}]
What does \ru{простуда} mean? (\textquotedblleft{}A cold\textquotedblright{}.) Say it once more.
\end{tcolorbox}
\section[\texorpdfstring{\ru{грипп} (gripp)}{грипп (gripp)}]{\ru{грипп} (gripp) — flu}

\label{lesson:RU-C147-gripp}



\begin{quote}
Before the new one: say the Russian for ill; a patient, then the Russian for a cold.
\end{quote}

\subsection*{You'll want to know: \ru{грипп}}
\textbf{\ru{грипп}} (\emph{gripp}) — \textquotedblleft{}flu\textquotedblright{}.

\begin{grammarlens}[title={What you've built}]
One more word for this chapter.
\end{grammarlens}

\subsection*{Your turn}
\begin{itemize}
\raggedright
  \item \emph{Say it:} \emph{gripp}
  \item \emph{Say it:} \emph{gripp}, once more
  \item \emph{Say it:} say \emph{gripp}
  \item \emph{Recall:} say the Russian for ill; a patient, then the Russian for a cold, then say \emph{gripp} again
\end{itemize}

\begin{tcolorbox}[breakable,colback=teal!4,colframe=teal!35!black,title={Before you move on}]
What does \ru{грипп} mean? (\textquotedblleft{}Flu\textquotedblright{}.) Say it once more.
\end{tcolorbox}
\section[\texorpdfstring{\ru{насморк} (násmork)}{насморк (násmork)}]{\ru{насморк} (násmork) — a runny nose}

\label{lesson:RU-C147-nasmork}



\begin{quote}
Before the new one: say the Russian for a cold, then the Russian for flu.
\end{quote}

\subsection*{You'll want to know: \ru{насморк}}
\textbf{\ru{насморк}} (\emph{násmork}) — \textquotedblleft{}a runny nose\textquotedblright{}.

\begin{grammarlens}[title={What you've built}]
One more word for this chapter.
\end{grammarlens}

\subsection*{Your turn}
\begin{itemize}
\raggedright
  \item \emph{Say it:} \emph{násmork}
  \item \emph{Say it:} \emph{násmork}, once more
  \item \emph{Say it:} say \emph{násmork}
  \item \emph{Recall:} say the Russian for a cold, then the Russian for flu, then say \emph{násmork} again
\end{itemize}

\begin{tcolorbox}[breakable,colback=teal!4,colframe=teal!35!black,title={Before you move on}]
What does \ru{насморк} mean? (\textquotedblleft{}A runny nose\textquotedblright{}.) Say it once more.
\end{tcolorbox}
\section[\texorpdfstring{\ru{таблетка} (tablétka)}{таблетка (tablétka)}]{\ru{таблетка} (tablétka) — a pill}

\label{lesson:RU-C147-tabletka}



\begin{quote}
Before the new one: say the Russian for flu, then the Russian for a runny nose.
\end{quote}

\subsection*{You'll want to know: \ru{таблетка}}
\textbf{\ru{таблетка}} (\emph{tablétka}) — \textquotedblleft{}a pill\textquotedblright{}.

\begin{grammarlens}[title={What you've built}]
One more word for this chapter.
\end{grammarlens}

\subsection*{Your turn}
\begin{itemize}
\raggedright
  \item \emph{Say it:} \emph{tablétka}
  \item \emph{Say it:} \emph{tablétka}, once more
  \item \emph{Say it:} say \emph{tablétka}
  \item \emph{Recall:} say the Russian for flu, then the Russian for a runny nose, then say \emph{tablétka} again
\end{itemize}

\begin{tcolorbox}[breakable,colback=teal!4,colframe=teal!35!black,title={Before you move on}]
What does \ru{таблетка} mean? (\textquotedblleft{}A pill\textquotedblright{}.) Say it once more.
\end{tcolorbox}
